%-------------------------
% ATS-Friendly Resume — Md Raihanul Haque Rahi
% Based on Jake Gutierrez's template (MIT License)
% Compiled with: pdflatex resume.tex
%-------------------------

\documentclass[a4paper,11.5pt]{extarticle}

\usepackage{latexsym}
\usepackage[a4paper,left=0.55in,right=0.55in,top=0.60in,bottom=0.60in]{geometry}
\usepackage{titlesec}
\usepackage[usenames,dvipsnames]{color}
\usepackage{verbatim}
\usepackage{enumitem}
\usepackage[hidelinks]{hyperref}
\usepackage{fancyhdr}
\usepackage[english]{babel}
\usepackage{tabularx}
\usepackage{microtype}
\usepackage{needspace}
\usepackage{ragged2e}
\usepackage{iftex}
\ifPDFTeX
  \input{glyphtounicode}
\fi

\pagestyle{fancy}
\fancyhf{}
\fancyfoot{}
\renewcommand{\headrulewidth}{0pt}
\renewcommand{\footrulewidth}{0pt}
\setlength{\footskip}{4pt}

\urlstyle{same}
\raggedbottom
\justifying
\hyphenpenalty=10000
\exhyphenpenalty=10000
\emergencystretch=2em
\setlength{\tabcolsep}{0in}

\titleformat{\section}{
  \vspace{2pt}\scshape\raggedright\large
}{}{0em}{}[\color{black}\titlerule\vspace{-2pt}]

\ifPDFTeX
  \pdfgentounicode=1
\fi

%-------- CUSTOM COMMANDS --------

\newcommand{\resumeItem}[1]{
  \item\normalsize{#1\vspace{1pt}}
}

\newcommand{\resumeSubheading}[4]{
  \vspace{2pt}\item
    \begin{tabular*}{0.97\textwidth}[t]{l@{\extracolsep{\fill}}r}
      \textbf{#1} & \small#2 \\
      \textit{#3} & \textit{\small#4} \\
    \end{tabular*}\vspace{-4pt}
}

\newcommand{\resumeProjectHeading}[2]{
  \vspace{2pt}\item
    \begin{tabular*}{0.97\textwidth}{l@{\extracolsep{\fill}}r}
      #1 & \small#2 \\
    \end{tabular*}\vspace{-4pt}
}

\renewcommand\labelitemii{$\vcenter{\hbox{\tiny$\bullet$}}$}

\newcommand{\resumeSubHeadingListStart}{\begin{itemize}[leftmargin=0.15in, label={}, topsep=2pt]}
\newcommand{\resumeSubHeadingListEnd}{\end{itemize}\vspace{2pt}}
\newcommand{\resumeItemListStart}{\begin{itemize}[leftmargin=0.2in, itemsep=1pt, topsep=1pt]}
\newcommand{\resumeItemListEnd}{\end{itemize}\vspace{2pt}}

%===========================================
\begin{document}

%-----------HEADER-----------
\begin{center}
  {\Huge\bfseries\scshape Md Raihanul Haque Rahi}\\[6pt]
  \small
  \href{mailto:raihanulhaque007@gmail.com}{raihanulhaque007@gmail.com}
  \enspace$\cdot$\enspace
  {+880\,1725-737812}
  \enspace$\cdot$\enspace
  \href{https://cortextech.dev/engineer/rahi/}{cortextech.dev/engineer/rahi}
  \enspace$\cdot$\enspace
  \href{https://github.com/raihanulhaque}{github.com/raihanulhaque}
  \enspace$\cdot$\enspace
  \href{https://linkedin.com/in/raihanulhaque}{linkedin.com/in/raihanulhaque}
\end{center}

\vspace{2pt}

%-----------SUMMARY-----------
\section{Summary}
\normalsize{
  AI Engineer with 2+ years of experience building and deploying LLM applications, including RAG chatbots,
  AI agents, and recommendation systems serving millions of users in fintech and media. Experienced in
  optimizing models for on-device inference on mobile and edge devices with llama.cpp, TFLite, and ONNX.
}

\vspace{2pt}

%-----------EDUCATION-----------
\section{Education}
\resumeSubHeadingListStart
  \resumeSubheading
    {Rajshahi University of Engineering \& Technology (RUET)}{Rajshahi, Bangladesh}
    {B.Sc.\ in Mechatronics Engineering $|$ CGPA: 3.16/4.00}{2022 -- 2026}
\resumeSubHeadingListEnd

%-----------TECHNICAL SKILLS-----------
\section{Technical Skills}
\begin{itemize}[leftmargin=0.15in, label={}, itemsep=3pt, topsep=2pt]
  \normalsize\item{
    \textbf{Languages:}\quad Python, C/C++, Dart, Rust, Swift, Kotlin, SQL \\
    \textbf{Frameworks \& Libraries:}\quad FastAPI, Django, Flutter, PyTorch, TensorFlow, MLX, LangChain, LiteLLM, HuggingFace, llama.cpp \\
    \textbf{Databases \& Infra:}\quad PostgreSQL, Supabase, Redis, Firebase, Docker, OpenRouter, WebRTC, LiveKit, Ray \\
    \textbf{AI / ML:}\quad Agentic LLM workflows, LLM routing, RAG, LLM fine-tuning (TRL, Unsloth, LoRA/QLoRA), ASR/TTS/VAD, YOLO, stereo depth, TFLite, ONNX, MLX, GGUF quantization
  }
\end{itemize}

\vspace{2pt}

%-----------EXPERIENCE-----------
\section{Experience}
\resumeSubHeadingListStart

  \resumeSubheading
    {Sysnova Information Systems}{Oct 2026 -- Present}
    {AI Engineer}{Dhaka, Bangladesh}
    \resumeItemListStart
      \resumeItem{Built a \textbf{RAG chatbot} for Prime Bank's mobile financial service \textbf{Lenden}, serving around \textbf{700k users}, and a \textbf{recommendation engine} for the Deepto Play OTT platform (\textbf{8.7M} total, \textbf{1M} active users)}
      \resumeItem{Developed an internal \textbf{AI coding gateway} on \textbf{LiteLLM} and \textbf{OpenRouter} that routes each request to cheap, mid, or strong models, with per-developer keys and budgets; A/B tested a hosted \textbf{JEV} decision model as the routing classifier}
      \resumeItem{Developed an \textbf{AI agent} for internal \textbf{LibreOffice Writer and Calc} that automates spreadsheet and document work}
    \resumeItemListEnd

  \resumeSubheading
    {Intelsense AI}{Apr 2024 -- Sep 2026}
    {Backend Engineer}{Remote}
    \resumeItemListStart
      \resumeItem{Owned the backend architecture for an AI-powered SaaS platform, designing \textbf{Django}/\textbf{PostgreSQL} systems with a focus on modularity, reliability, and scale}
      \resumeItem{Architected production \textbf{RAG} services with structured ingestion and retrieval flows, resilient failure handling, and reliable product-facing APIs}
      \resumeItem{Built \textbf{OCR and ASR data-collection pipelines} and integrated LLM, speech, and document-processing capabilities into the web backend}
    \resumeItemListEnd

  \resumeSubheading
    {Cortex Technologies}{Jan 2024 -- Present}
    {ML Engineer}{Freelance Collaboration}
    \resumeItemListStart
      \resumeItem{Own the ML engineering track for client projects sourced through a small freelance collaboration, building custom object detection pipelines, stereo depth estimation, and real-time edge-vision inference}
      \resumeItem{Built an automated object-detection dataset pipeline that uses \textbf{SAM 3} to generate bounding boxes from user-provided labels and exports annotations in \textbf{COCO format}}
      \resumeItem{Optimize and deploy models on resource-constrained hardware using \textbf{TFLite}, \textbf{ONNX}, and quantization techniques}
    \resumeItemListEnd

  \resumeSubheading
    {Outlier / Remotasks}{Jul 2023 -- Dec 2024}
    {AI Trainer \& LLM Data Quality Reviewer}{Remote}
    \resumeItemListStart
      \resumeItem{Reviewed AI-generated math/code solutions and audited multilingual programming datasets for model reasoning and code-generation fine-tuning quality}
    \resumeItemListEnd

\resumeSubHeadingListEnd

%-----------PROJECTS-----------
\section{Projects}
\resumeSubHeadingListStart

  \resumeProjectHeading
    {\textbf{\href{https://github.com/RaihanulHaque/kiln}{Kiln: On-Device LLM Runtime}} $|$ \emph{C/C++, llama.cpp, Flutter, Kotlin, Rust, React Native}}{Oct 2026 -- Present}
    \resumeItemListStart
      \resumeItem{Built an open-source \textbf{on-device LLM runtime} on llama.cpp for Android and desktop, serving \textbf{chat, decision, embedding, speech-to-text, and text-to-speech} models through one C API, with bindings for Flutter (published on pub.dev as \texttt{kiln\_flutter}), Kotlin, Rust, and React Native}
      \resumeItem{Detects hardware acceleration at runtime and falls back to a plain Armv8 CPU build on older phones; reports the loaded backend, memory use, and real tokens/s; verified with \textbf{505 C checks} across Android, macOS, and Linux}
    \resumeItemListEnd

  \resumeProjectHeading
    {\textbf{StereoLite: Lightweight Stereo Depth Estimation} $|$ \emph{PyTorch, ONNX $\cdot$ B.Sc.\ Thesis}}{Active}
    \resumeItemListStart
      \resumeItem{Designed a stereo depth model with \textbf{fewer than 3M parameters}, over \textbf{10$\times$ smaller} than DEFOM-Stereo (30M+), while using \textbf{4$\times$ less VRAM} (1.5 GB versus 6 GB) and reaching EPE \textbf{0.78} on the full SceneFlow benchmark}
      \resumeItem{Trained purely on synthetic SceneFlow data, yet still generalizes zero-shot to real-world Middlebury 2014 scenes with a D1 of 10.86\%, showing the model holds up outside its training distribution}
    \resumeItemListEnd

  \resumeProjectHeading
    {\textbf{AgriSense: Autonomous Agricultural Advisor} $|$ \emph{LangGraph, RAG, FastAPI}}{IUT Hackathon}
    \resumeItemListStart
      \resumeItem{Built a \textbf{multi-agent system} with LangGraph that plans a farmer's growing season, including costs and weather, and keeps giving advice as the season goes on}
      \resumeItem{Used real government data on soil, weather, and crop needs by district and upazila, plus a crop leaf-disease classifier, so the advice matches what a real soil scientist would recommend; deployed with CI/CD at \href{https://agrisense.cortextech.dev/}{\texttt{agrisense.cortextech.dev}}}
    \resumeItemListEnd

  \resumeProjectHeading
    {\textbf{Hishab: Offline AI Expense Tracker} $|$ \emph{Python, PyTorch, Qwen3.5, GGUF, Kiln, Flutter}}{May 2024 -- Present}
    \resumeItemListStart
      \resumeItem{Fine-tuned \textbf{Qwen3.5-0.8B} to turn English, Bangla, Banglish, and mixed-language messages into structured income and expense records, shipped as an offline Android app: the \textbf{343 MB} quantized model runs on the phone through Kiln at about \textbf{51 tokens/s} on a Pixel 8, with no server and no account}
      \resumeItem{Relabelled the synthetic training data against a written guideline and trimmed the vocabulary from 248k to 38k tokens, shrinking the model from 529 MB to 343 MB and speeding up generation by about 45\%, while reaching \textbf{94.2\%} exact-match on a cleaned held-out test set}
    \resumeItemListEnd

  \resumeProjectHeading
    {\textbf{Gesprek} $|$ \emph{FastAPI, PostgreSQL, Redis, Supabase, Docker, ElevenLabs, LiveKit}}{Archived}
    \resumeItemListStart
      \resumeItem{Built and deployed a \textbf{production-ready voice-agent platform} with configurable multi-agent workflows, Supabase auth, Redis-backed real-time TTS, subscriptions, and containerized services}
      \resumeItem{Engineered multi-tenant agent orchestration and independently scalable services; archived after launch since the project could not secure enough funding to take it to market}
    \resumeItemListEnd

\resumeSubHeadingListEnd

%-----------PUBLICATIONS-----------
\section{Publication}
\begin{itemize}[leftmargin=0.2in, itemsep=1pt, topsep=1pt]
  \normalsize{
    \item \textbf{``Diet Engine: A real-time food nutrition assistant system for personalized dietary guidance.''}
          \textit{Food Chemistry Advances}, Elsevier, 2025. \href{https://doi.org/10.1016/j.focha.2025.100978}{doi:10.1016/j.focha.2025.100978}
  }
\end{itemize}

%-----------ACHIEVEMENTS-----------
\section{Achievements \& Activities}
\begin{itemize}[leftmargin=0.2in, itemsep=1pt, topsep=1pt]
  \normalsize{
    \item \textbf{JAXA Kibo Robot Programming Challenge 2025:} \textbf{1st place, Bangladesh regional}; \textbf{4th place globally};
          programmed Astrobee robots for zero-gravity navigation and AR object recognition aboard the ISS
    \item \textbf{Microsoft / VibeVoice:} Contributed package-import support to a 40k+ star emotional voice-synthesis repository
    \item \textbf{Leadership:} Codeforces rating \textbf{1280} (Pupil); led RUET coding club and taught C/C++ problem-solving to first-year students; runs Rahi 404 technical YouTube channel
  }
\end{itemize}

\end{document}
